\centerline{\bf \Huge Банки. База(А+Д) II. 2}
\centerline{\bf \Huge }

\begin{flushright}
        \scriptsize{\textbf{qq}}
\end{flushright}


\problem{\textbf{1}} В июле 2025 года планируется взять кредит в банке на некоторую сумму на 10 лет. Условия его возврата таковы:

– каждый январь долг будет возрастать на 30\% по сравнению с концом предыдущего года;

– с февраля по июнь каждого года необходимо выплатить одним платежом часть долга;

– в июле 2026, 2027, 2028, 2029 и 2030 годов долг должен быть на какую-то одну и ту же
величину меньше долга на июль предыдущего года;

– в июле 2030 года долг должен составлять 700 тыс. рублей;

– в июле 2031, 2032, 2033, 2034 и 2035 годов долг должен быть на другую одну и ту же
величину меньше долга на июль предыдущего года;

– к июлю 2035 года кредит должен быть выплачен полностью.

Известно, что сумма всех платежей после полного погашения кредита будет равна 2760 тыс.
рублей. Какую сумму планируется взять в кредит?

\problem{\textbf{2}} В июле 2025 года планируется взять кредит на 600 тыс. рублей. Условия его возврата таковы:

– в январе 2026, 2027, 2028, 2029 и 2030 годов долг возрастает на 13\% по сравнению с концом
предыдущего года;

– в январе 2031, 2032, 2033, 2034, 2035 годов долг возрастает на 12\% по сравнению с концом
предыдущего года;

– в июле каждого года долг должен быть на одну и ту же величину меньше долга на июль
предыдущего года;

– к июлю 2035 года долг должен быть полностью погашен.

Чему равна сумма всех выплат?


\problem{\textbf{3}} 15-го декабря планируется взять кредит в банке на 1000 тыс. рублей на 21 месяц. Условия его возврата таковы:

– 1-го числа каждого месяца долг возрастает на $r$ \% по сравнению с концом предыдущего
месяца;

– со 2-го по 14-е число каждого месяца необходимо выплатить часть долга;

– 15-го числа с 1-го по 20-й месяцы долг должен быть на одну и ту же сумму меньше долга
на 15-е число предыдущего месяца.

Известно, что общая сумма выплат после полного погашения кредита составила 1441 тыс.
рублей, а долг на 15-е число 20-го месяца — 400 тыс. рублей. Найдите $r$